% !TEX root = main.tex
\clearpage
\section{Lesson 8 --- Training Engineering}

Building a Transformer defines what the model can compute. Training engineering
defines whether that model can learn reliably, efficiently, and reproducibly.
This lesson turns the equations from the previous lesson into a complete loop
that can run on a CPU, a CUDA GPU, or Apple Silicon through MPS.

\subsection*{The complete training loop}

\begin{center}
\begin{tikzpicture}[
  node distance=0.28cm and 0.8cm,
  >=Latex,
  font=\scriptsize,
  box/.style={draw, rounded corners, align=center, minimum height=0.55cm,
              minimum width=4.2cm, inner xsep=0.35cm, fill=gray!5},
  side/.style={box, minimum width=3.1cm},
  flow/.style={->, semithick}
]
  \node[box] (dataset) {Dataset};
  \node[box, below=of dataset] (batch) {Create batch\\inputs + targets};
  \node[box, below=of batch] (device) {Move batch to device\\CPU / CUDA GPU / MPS};
  \node[box, below=of device] (precision) {Mixed precision\\FP32 / BF16 / FP16};
  \node[box, below=of precision] (forward) {MiniGPT\\forward pass};
  \node[box, below=of forward] (loss) {Cross-entropy loss};
  \node[box, below=of loss] (backward) {Backpropagation\\\texttt{loss.backward()}};
  \node[box, below=of backward] (grads) {Gradients are created};
  \node[side, below left=0.35cm and 0.2cm of grads] (accum) {Gradient accumulation\\``combine batches''};
  \node[side, below right=0.35cm and 0.2cm of grads] (clip) {Gradient clipping\\``stop huge gradients''};
  \node[box, below=1.35cm of grads] (schedule) {Learning-rate schedule\\warmup $\longrightarrow$ decay};
  \node[box, below=of schedule] (adamw) {AdamW\\update model weights};
  \node[side, below left=0.35cm and 0.2cm of adamw] (evaluate) {Evaluate model\\train / validation loss\\perplexity};
  \node[side, below right=0.35cm and 0.2cm of adamw] (checkpoint) {Save checkpoint\\model + optimizer + step};
  \node[box, below=1.55cm of adamw] (next) {Next training step};

  \draw[flow] (dataset) -- (batch);
  \draw[flow] (batch) -- (device);
  \draw[flow] (device) -- (precision);
  \draw[flow] (precision) -- (forward);
  \draw[flow] (forward) -- (loss);
  \draw[flow] (loss) -- (backward);
  \draw[flow] (backward) -- (grads);
  \draw[flow] (grads) -| (accum);
  \draw[flow] (grads) -| (clip);
  \draw[flow] (accum) |- (schedule);
  \draw[flow] (clip) |- (schedule);
  \draw[flow] (schedule) -- (adamw);
  \draw[flow] (adamw) -| (evaluate);
  \draw[flow] (adamw) -| (checkpoint);
  \draw[flow] (evaluate) |- (next);
  \draw[flow] (checkpoint) |- (next);
  \draw[flow] (next.east) -- ++(0.55,0) |- (dataset.east);
\end{tikzpicture}
\end{center}

A \emph{micro-step} processes one mini-batch and creates gradients. An
\emph{optimizer step} changes the weights. They are the same only when
gradient accumulation is one.

\subsection{Dataset and next-token batches}

The dataset is a stream of token IDs. For sequence length (T), each example
needs (T+1) consecutive tokens. The first (T) are inputs and the last
(T), shifted left by one position, are targets:

\[
  x=(t_0,t_1,\ldots,t_{T-1}),
  \qquad
  y=(t_1,t_2,\ldots,t_T).
\]

For batch size (B), both tensors have shape ([B,T]).

\begin{lstlisting}[style=shadedpython]
def get_batch(token_ids, batch_size, sequence_length, device):
    starts = torch.randint(
        0, len(token_ids) - sequence_length - 1, (batch_size,)
    )
    x = torch.stack([
        token_ids[i : i + sequence_length] for i in starts
    ])
    y = torch.stack([
        token_ids[i + 1 : i + sequence_length + 1] for i in starts
    ])
    return x.to(device), y.to(device)
\end{lstlisting}

Training and validation data must be separate. If validation examples leak
into training, validation loss no longer measures generalization.

\subsection{Move the batch to the compute device}

\begin{lstlisting}[style=shadedpython]
def get_device():
    if torch.cuda.is_available():
        return torch.device("cuda")

    if torch.backends.mps.is_available():
        return torch.device("mps")

    return torch.device("cpu")


device = get_device()

print("Using:", device)
\end{lstlisting}

\subsubsection*{Moving the model onto the GPU}

Creating

\begin{lstlisting}[style=shadedpython]
model = MiniGPT()
\end{lstlisting}

initially places its parameters on the CPU. Move everything with

\begin{lstlisting}[style=shadedpython]
model = model.to(device)
\end{lstlisting}

Now the following components live in accelerator memory:

\begin{itemize}
  \setlength{\itemsep}{0pt}
  \item token embeddings,
  \item Q/K/V matrices,
  \item feed-forward weights,
  \item LayerNorm parameters, and
  \item the output projection.
\end{itemize}

\subsubsection*{Inputs must be on the same device}

This configuration will fail:

\[
  \text{Model}\longrightarrow\text{GPU},
  \qquad
  \text{Input}\longrightarrow\text{CPU}.
\]

The batch generator should return tensors on the same device as the model. \\

Instead of

\begin{lstlisting}[style=shadedpython]
return inputs, targets
\end{lstlisting}

use

\begin{lstlisting}[style=shadedpython]
return (
    inputs.to(device),
    targets.to(device),
)
\end{lstlisting}

Now the placement is consistent:

\begin{lstlisting}[style=shadedpython,language={}]
Model   -> GPU
Inputs  -> GPU
Targets -> GPU
\end{lstlisting}

GPU memory contains roughly

\[
  \begin{aligned}
  &\text{model parameters}
  +\text{gradients}
  +\text{optimizer states}\\
  &\quad+\text{activations}
  +\text{attention matrices}
  +\text{temporary computation buffers}.
  \end{aligned}
\]

This becomes very important when scaling.

\subsection{Mixed precision: FP32, BF16, and FP16}

FP32 is stable but consumes more memory. Mixed precision performs expensive
operations in a lower-precision format while keeping sensitive state, such as
optimizer statistics, in FP32.

\begin{center}
\renewcommand{\arraystretch}{1.2}
\begin{tabular}{@{}p{0.14\textwidth}p{0.27\textwidth}p{0.47\textwidth}@{}}
  \textbf{Format} & \textbf{Main advantage} & \textbf{Practical use} \\
  \hline
  FP32 & Numerical range and precision & Debugging and hardware without reliable mixed precision \\
  BF16 & FP32-like exponent range & Preferred on modern accelerators when supported \\
  FP16 & High throughput and low memory & Often needs dynamic loss scaling to avoid underflow
\end{tabular}
\end{center}

\begin{lstlisting}[style=shadedpython]
amp_enabled = device.type == "cuda"
amp_dtype = torch.bfloat16 if bf16_supported else torch.float16

with torch.autocast(
    device_type=device.type,
    dtype=amp_dtype,
    enabled=amp_enabled,
):
    logits, loss = model(x, y)
\end{lstlisting}

For FP16, a gradient scaler multiplies the loss before backpropagation and
unscales gradients before clipping. BF16 usually does not need loss scaling
because it has a much wider exponent range.

\subsection{Forward pass and cross-entropy loss}

The forward pass maps input tokens to logits:

\[
  [B,T]\longrightarrow\operatorname{MiniGPT}(x)
  \longrightarrow Z\in\mathbb{R}^{B\times T\times V}.
\]

Cross-entropy compares every row of vocabulary logits with the correct next
token. With \(N=B\times T\) predictions, the mean loss is

\[
  L=-\frac{1}{N}\sum_{i=1}^{N}\log p(y_i).
\]

This scalar starts backpropagation. A finite loss does not by itself prove
that training is healthy; it must trend downward on training data and remain
meaningful on held-out validation data.

\subsection{Backpropagation creates gradients}

Before a new optimizer step, stale gradients must be cleared. Then
\texttt{loss.backward()} applies the chain rule and adds a gradient to every
trainable parameter's \texttt{.grad} field.

\begin{lstlisting}[style=shadedpython]
optimizer.zero_grad(set_to_none=True)
logits, loss = model(x, y)
loss.backward()
\end{lstlisting}

Backpropagation does not update the weights. It computes

\[
  g_t=\nabla_{\theta}L_t.
\]

PyTorch accumulates gradients by addition. This is why gradients must be
cleared deliberately, and what makes gradient accumulation possible.

\subsection{Gradient accumulation}

When the desired batch is too large for device memory, process (K) smaller
micro-batches, divide each loss by (K), and call backward on each one before
a single optimizer update:

\[
  B_{\text{effective}}
  =B_{\text{micro}}\times K\times N_{\text{workers}}.
\]

\begin{lstlisting}[style=shadedpython]
optimizer.zero_grad(set_to_none=True)

for _ in range(accumulation_steps):
    x, y = get_batch(...)
    logits, loss = model(x, y)
    (loss / accumulation_steps).backward()

optimizer.step()
\end{lstlisting}

Dividing by (K) preserves the scale of the mean gradient. The learning-rate
schedule and global step normally advance once per optimizer step, not once
per micro-batch.

\subsection{FP16 and gradient scaling}

FP16 can represent only a limited range of very small numbers.

During backpropagation, some gradients may become tiny:

\[
  0.00000000013
\]

FP16 might effectively round such a value down to:

\[
  0
\]

Then learning stops for that gradient.

One solution is \textbf{gradient scaling}.

Instead of backpropagating gradients directly, we temporarily use:

\[
  S \times L
\]

where $S$ is a large scaling factor.

This makes gradients larger. Afterward, they are scaled back down before updating the parameters.

PyTorch's:

\begin{lstlisting}[style=shadedpython]
torch.amp.GradScaler
\end{lstlisting}

automates this for FP16 CUDA training.

BF16 usually does not require gradient scaling because its exponent range is much larger.

\subsection{Why gradients can explode}

Imagine your model has many layers.

Backpropagation repeatedly multiplies derivatives through them.

Occasionally gradients can become huge:

\[
  \text{normal gradient: } 0.04
\]

\[
  \text{problematic gradient: } 9482
\]

Then one optimizer step might dramatically modify the weights.

You may see:

\begin{lstlisting}[style=shadedpython,language={}]
loss = 2.8
loss = 2.4
loss = 2.1
loss = NaN
\end{lstlisting}

Training has become unstable.

\subsection{Gradient clipping}

We can constrain the total gradient norm.

In PyTorch:

\begin{lstlisting}[style=shadedpython]
torch.nn.utils.clip_grad_norm_(
    model.parameters(),
    max_norm=1.0,
)
\end{lstlisting}

This does \textbf{not} simply clamp every gradient between $-1$ and $1$.

Instead, it calculates an overall gradient norm.

Suppose:

\begin{lstlisting}[style=shadedpython,language={}]
gradient norm = 8
maximum = 1
\end{lstlisting}

PyTorch scales the gradients down proportionally.

Their directions remain roughly the same.

\subsection{Where gradient clipping goes}

Training should now look like:

\begin{lstlisting}[style=shadedpython]
optimizer.zero_grad(set_to_none=True)

logits, loss = model(inputs, targets)

loss.backward()

torch.nn.utils.clip_grad_norm_(
    model.parameters(),
    1.0,
)

optimizer.step()
\end{lstlisting}

The order matters:

\[
\begin{aligned}
&\text{forward} \\
&\downarrow \\
&\text{calculate loss} \\
&\downarrow \\
&\text{backward} \\
&\downarrow \\
&\text{gradients exist} \\
&\downarrow \\
&\text{clip gradients} \\
&\downarrow \\
&\text{optimizer updates weights}
\end{aligned}
\]

\subsection{Constant learning rate is not ideal}

Until now we used:

\begin{lstlisting}[style=shadedpython,language={}]
learning_rate = 3e-4
\end{lstlisting}

for every training step.

That means:

\begin{lstlisting}[style=shadedpython,language={}]
Step 1      → 0.0003
Step 100    → 0.0003
Step 10,000 → 0.0003
Step 50,000 → 0.0003
\end{lstlisting}

Real training often changes the learning rate over time.

A common schedule is:

\[
\begin{aligned}
&\text{warmup} \\
&\downarrow \\
&\text{maximum learning rate} \\
&\downarrow \\
&\text{gradual decay}
\end{aligned}
\]

\subsection{Why warmup?}

At initialization, the model is essentially random.

Gradients can be noisy.

If we immediately use a large learning rate:

\[
\begin{aligned}
&\text{random model} \\
&+ \\
&\text{large updates} \\
&= \\
&\text{potential instability}
\end{aligned}
\]

Instead, we might start at:

\[
  0.000003
\]

and gradually rise to:

\[
  0.0003
\]

over the first few hundred or thousand steps.

This is \textbf{learning-rate warmup}.

\subsection{Warmup calculation}

Suppose:

\begin{lstlisting}[style=shadedpython,language={}]
warmup_steps = 1000
max_learning_rate = 3e-4
\end{lstlisting}

At step 500:

\[
  \text{LR} = \frac{500}{1000} \times 0.0003 = 0.00015
\]

At step 1000:

\[
  \text{LR} = 0.0003
\]

\subsection{Learning-rate decay}

After warmup, we reduce the learning rate.

Early training:

\begin{lstlisting}[style=shadedpython,language={}]
make relatively large improvements
\end{lstlisting}

Later training:

\begin{lstlisting}[style=shadedpython,language={}]
make small refinements
\end{lstlisting}

A common schedule is cosine decay.

Conceptually:

\[
\begin{array}{c}
\text{LR} \\
\\
\mid \quad \quad /\backslash \\
\mid \quad /\quad\backslash \\
\mid \quad /\quad\quad\backslash \\
\mid \quad /\quad\quad\quad\backslash \\
\mid \quad /\quad\quad\quad\quad\backslash \\
\mid \_\_\_/\quad\quad\quad\quad\quad\backslash\_\_\_ \\
\mid \\
\text{warmup} \quad\quad\quad \text{decay}
\end{array}
\]

\subsection{A learning-rate function}

Here's a practical implementation:

\begin{lstlisting}[style=shadedpython]
import math

def get_learning_rate(step):
    warmup_steps = 500
    max_steps = training_steps

    max_lr = 3e-4
    min_lr = 3e-5

    # Warmup
    if step < warmup_steps:
        return (
            max_lr
            * (step + 1)
            / warmup_steps
        )

    # Finished decay
    if step >= max_steps:
        return min_lr

    # Cosine decay progress: 0 -> 1
    progress = (
        (step - warmup_steps)
        /
        (max_steps - warmup_steps)
    )

    cosine = 0.5 * (
        1.0
        +
        math.cos(
            math.pi * progress
        )
    )

    return (
        min_lr
        +
        cosine
        * (max_lr - min_lr)
    )
\end{lstlisting}

\subsection{AdamW updates the weights}

Once gradients are ready, AdamW updates its first and second moments and then
changes the parameters. In simplified form,

\[
  m_t=\beta_1m_{t-1}+(1-\beta_1)g_t,
  \qquad
  v_t=\beta_2v_{t-1}+(1-\beta_2)g_t^2,
\]
\[
  \theta_t=\theta_{t-1}
  -\eta_t\frac{\widehat m_t}{\sqrt{\widehat v_t}+\epsilon}
  -\eta_t\lambda\theta_{t-1}.
\]

The final term is decoupled weight decay. Biases and normalization parameters
are commonly assigned zero weight decay.

\begin{lstlisting}[style=shadedpython]
optimizer = torch.optim.AdamW(
    parameter_groups,
    lr=max_lr,
    betas=(0.9, 0.95),
    eps=1e-8,
)
\end{lstlisting}

Hyperparameters must be recorded with the code version, data version, random
seed, and model configuration.

\subsection{Evaluation and perplexity}

Evaluation measures performance without creating gradients or changing
weights. Switch to evaluation mode, sample fixed validation batches, then
restore training mode.

\begin{lstlisting}[style=shadedpython]
@torch.no_grad()
def estimate_loss(model, get_validation_batch, eval_steps):
    model.eval()
    losses = []
    for _ in range(eval_steps):
        x, y = get_validation_batch()
        _, loss = model(x, y)
        losses.append(loss.detach().float())
    model.train()
    return torch.stack(losses).mean().item()
\end{lstlisting}

For token cross-entropy measured with natural logarithms,

\[
  \operatorname{PPL}=e^{L_{\text{validation}}}.
\]

Perplexity is comparable only when tokenization, evaluation data, and loss
conventions match. Task-specific evaluations are still needed; loss alone
does not measure factuality, coding, or financial reasoning.

\subsection{Save resumable checkpoints}

A checkpoint needs enough state to resume the same optimization trajectory:

\begin{itemize}
  \setlength{\itemsep}{0pt}
  \item model parameters;
  \item optimizer, scheduler, and FP16 gradient-scaler state;
  \item optimizer step and consumed-token count;
  \item model and training configuration;
  \item random-number-generator states; and
  \item data sampler position when deterministic resumption matters.
\end{itemize}

\begin{lstlisting}[style=shadedpython]
checkpoint = {
    "model": model.state_dict(),
    "optimizer": optimizer.state_dict(),
    "scaler": scaler.state_dict() if scaler is not None else None,
    "step": step,
    "tokens_seen": tokens_seen,
    "model_config": model_config,
    "train_config": train_config,
    "rng_state": torch.get_rng_state(),
}
torch.save(checkpoint, checkpoint_path)
\end{lstlisting}

Write to a temporary path and rename only after the write succeeds. Keep a
recent rolling checkpoint for recovery and less frequent milestone
checkpoints for comparison and auditing.

\subsection{One complete optimizer step}

The pieces must occur in the right order:

\begin{lstlisting}[style=shadedpython]
model.train()

for step in range(start_step, total_steps):
    lr = learning_rate(
        step, warmup_steps, total_steps, max_lr, min_lr
    )
    for group in optimizer.param_groups:
        group["lr"] = lr

    optimizer.zero_grad(set_to_none=True)
    accumulated_loss = 0.0

    for _ in range(accumulation_steps):
        x, y = get_train_batch()
        with torch.autocast(
            device_type=device.type,
            dtype=amp_dtype,
            enabled=amp_enabled,
        ):
            _, loss = model(x, y)
            scaled_loss = loss / accumulation_steps

        if scaler is None:
            scaled_loss.backward()
        else:
            scaler.scale(scaled_loss).backward()
        accumulated_loss += loss.detach().float().item()

    if scaler is not None:
        scaler.unscale_(optimizer)

    grad_norm = torch.nn.utils.clip_grad_norm_(
        model.parameters(), max_norm=1.0
    )

    if scaler is None:
        optimizer.step()
    else:
        scaler.step(optimizer)
        scaler.update()

    if step % eval_interval == 0:
        validation_loss = estimate_loss(...)
        validation_ppl = math.exp(validation_loss)

    if step % checkpoint_interval == 0:
        save_checkpoint(...)
\end{lstlisting}

\[
  \boxed{
  \text{zero gradients}
  \rightarrow\text{forward}
  \rightarrow\text{loss}
  \rightarrow\text{backward}
  \rightarrow\text{unscale}
  \rightarrow\text{clip}
  \rightarrow\text{AdamW}
  }
\]

Evaluation and checkpointing happen at controlled intervals; then the next
optimizer step begins.

\subsection{Signals to monitor}

At minimum, record optimizer step, tokens processed, train loss, validation
loss, learning rate, gradient norm, step time, tokens per second, memory use,
and skipped FP16 updates. Investigate non-finite loss or gradients, persistent
validation degradation, unexpected throughput changes, or checkpoint resumes
that do not reproduce the expected next steps.

The loop succeeds only when it is both a learning system and a measurement
system: it changes the weights and leaves enough evidence to explain exactly
how those weights were produced.
